\section{Methodology}

\subsection*{3.1 Research Design}
This study employs a Qualitative Single Instrumental Case Study design. This approach is ideal for understanding the complex legal, technical, and administrative intersections required to draft a new municipal policy.

\subsection*{3.2 Research Locale}
The study is situated in the Municipality of Maigo, Lanao del Norte, a locale characterized by a strong presence of technical-vocational institutions (MSU-MCEST and TESDA) and a cluster of struggling commercial internet cafés.

\subsection*{3.3 Key Informants}
The researcher will utilize Purposive Sampling across three sectors:
\begin{itemize}
    \item \textbf{Governance:} LGU Officials (Mayor, SB Committee on Education, PESO Manager).
    \item \textbf{Educational:} Administrators from MSU-MCEST and TESDA Provincial/Municipal coordinators.
    \item \textbf{Private:} 3–5 registered internet café owners in Maigo.
\end{itemize}

\subsection*{3.4 Research Instruments}
Data will be collected via Semi-Structured Interview Guides and Documentary Analysis of local ordinances and TESDA facility standards.

\subsection*{3.5 Data Gathering and Analysis}
Data will be gathered through clearances, document retrieval, and interviews. The study will use Thematic Analysis and Triangulation to cross-reference business needs, legal boundaries, and educational standards to draft the final policy output.

\subsection*{3.6 Technical System Architecture (The ``Dual-Use'' Mechanism)}
For this policy to be accepted, the researcher must demonstrate how the commercial and educational environments are separated. The study proposes the utilization of Diskless (PXE Boot) Architecture, which is already common in high-end Philippine internet cafes.
\begin{itemize}
    \item \textbf{Environmental Segregation:} The study will explore how a single server can host two distinct ``Images'':
    \begin{itemize}
        \item \textit{The ``Edu-Image'' (Daytime):} A restricted OS containing only Microsoft Office, AutoCAD, Programming IDEs, and TESDA assessment tools. All gaming ports and social media are blocked at the BIOS/Network level.
        \item \textit{The ``Gamer-Image'' (Nighttime):} The standard commercial OS with high-fidelity games and unrestricted internet access.
    \end{itemize}
    \item \textbf{Data Volatility:} The research will investigate the ``Reboot-to-Restore'' feature, ensuring that any files or potential malware introduced during one session are completely wiped before the next session begins.
\end{itemize}

\subsection*{3.7 Proposed Implementation Roadmap (The Policy Lifecycle)}
The study will outline a four-phase transition for the Municipality of Maigo to adopt this framework:

\begin{table}[h!]
\centering
\begin{tabular}{|l|p{7cm}|p{5cm}|}
\hline
\textbf{Phase} & \textbf{Activity} & \textbf{Responsible Entities} \\ \hline
I. Vetting & Technical and physical audit of local cafes based on TESDA standards. & LGU, TESDA, MSU-MCEST \\ \hline
II. Accreditation & Granting of ``Certified Dual-Use Tech Space'' status to eligible MSMEs. & Maigo Municipal Council \\ \hline
III. Integration & Deployment of the ``Edu-Image'' software and instructor training. & University IT Dept / TESDA \\ \hline
IV. Operation & Commencement of daytime classes and monthly lease disbursement. & LGU Finance / PESO \\ \hline
\end{tabular}
\end{table}

\subsection*{3.10 Anticipated Policy Outcomes}
\begin{itemize}
    \item \textbf{Economic Resilience:} A 30–40\% increase in daytime revenue for participating MSMEs, preventing local business closures.
    \item \textbf{Educational Expansion:} A 50\% increase in the capacity of MSU-MCEST and TESDA to accept ICT scholars without new building construction.
    \item \textbf{Governance Innovation:} A scalable ``Maigo Model'' of Public-Private Partnership that can be replicated by other 4th and 5th-class municipalities in Lanao del Norte.
\end{itemize}